\documentclass[11pt,fleqn]{article}
\usepackage[margin=0.8in]{geometry}
\usepackage{amsmath,amssymb}
\usepackage[hidelinks]{hyperref}
\setlength{\mathindent}{0pt}

\newcommand{\X}{\mathcal X}
\newcommand{\C}{\mathcal C}
\newcommand{\Hy}{\mathcal H}
\newcommand{\err}{\operatorname{error}}
\newcommand{\eps}{\epsilon}

\begin{document}
\raggedright
\setlength{\parskip}{0.6em}

{\LARGE\bfseries Page 3: Sample Complexity vs. Computational Complexity\par}
\medskip

\section{What Does it Mean to ``Fit a Sample''}

We will say that a hypothesis is \textbf{consistent} if it matches the labels of all the examples in the training sample. We will say that a hypothesis's \textbf{empirical risk} is the fraction of training examples that it incorrectly classifies. The strategy of \textbf{Empirical Risk Minimization} (ERM) is the one in which a learner returns the $h \in \Hy$ that has the smallest empirical risk. When the setting is \textit{realizable}, this means that the smallest possible empirical risk is zero.

\section{The Finite-Class Bound}

This is a very important bound in PAC, and its proof is not super complicated.\par

Let $\Hy$ be some fixed, finite hypothesis class. Let $f$ be the target function, and let $D$ be the input distribution. We will draw $m$ independent samples from $D$ to act as our training sample $S$, and we will label them using $f$. If $\boxed{m \geq \left\lceil \frac{\ln |\Hy| + \ln(1/\delta)}{\epsilon} \right\rceil}$, then with probability at least $1 - \delta$ over $S$, every consistent hypothesis in $\Hy$ has a true generalization error at most $\eps$.

In the realizable setting, this means that a strategy of ERM can use this sample size to meet the PAC guarantee~\cite{kv,ssbd}.

Let's now go through the proof.

\subsection*{Lemma: $1 + x \leq e^x$}

This is a standard mathematical fact that will be important for later. Note that using Taylor's theorem with the Lagrange remainder $e^x = 1 + x + \frac{e^c}{2}x^2$ for some $c$ between $0$ and $x$. Thus, $1 + x = e^x - \frac{e^c}{2}x^2$, and since $\frac{e^c}{2}x^2 \geq 0$, this means that $1 + x \leq e^x$. Note that this is a fairly good approximation when $x$ is close to 0.

\subsection*{Step 1: Bad Rules}

For the first step in our proof, let us define a \textit{bad} rule as an $h$ where $\err_D(h, f) > \epsilon$, where $0 < \eps \leq 1$. Note that whether or not $h$ is bad is determined by $D$ and $f$. The randomness comes from whether or not the sample we draw will expose the mistake.

\subsection*{Step 2: Look at One Particular Bad Rule}

Let us say a fixed bad rule $h$ has an a true error $p > \epsilon$. Then, a single draw will miss its mistakes with probability $1 - p$. Thus, over $m$ i.i.d. draws, the probability that our bad rule never sees a mistake (so the probability that $h$ is consistent with $S$) is
\[(1 - p)^m \leq (1 - \eps)^m \leq e^{-\eps m}\]

\subsection*{Step 3: Account for All the Rules}

Note that there could be at most $|\Hy|$ bad hypotheses. Thus, using the union bound, we can say that
\[\boxed{\Pr_S [\text{at least one bad hypothesis is consistent with } S] \leq |\Hy|e^{-\eps m}}\]

\subsection*{Step 4: Solve for the Sample Size}

Now, note that we want $|\Hy|e^{-\eps m} \leq \delta$, where $0 < \delta \leq 1$. Thus,
\begin{align*}
    |\Hy|e^{-\eps m} \leq \delta &\iff e^{-\eps m} \leq \frac{\delta}{|\Hy|}\\
    &\iff -\eps m \leq \ln (\delta / |\Hy|)\\
    &\iff -\eps m \leq \ln (\delta) - \ln(|\Hy|)\\
    &\iff \eps m \geq \ln(|\Hy|) - \ln (\delta) \\
    &\iff m \geq \frac{\ln(|\Hy|) - \ln (\delta)}{\eps}\\
    &\iff m \geq \frac{\ln(|\Hy|) + \ln (1/\delta)}{\eps}
\end{align*}
Taking the ceiling so that we get a whole number, this leaves us with $\boxed{m \geq \left\lceil \frac{\ln |\Hy| + \ln(1/\delta)}{\epsilon} \right\rceil}$.


\section{Sufficient Data and the Cost of Searching}

Note that the general bound we just derived gives a sufficient number, but a specific problem may need fewer examples. For instance, for a hypothesis class of two rules with $\epsilon = \delta = 0.05$, the general bound gives $m \geq \lceil \ln (40) / 0.05\rceil = 74$ examples. However, let's say that exactly one of the rules is bad and makes mistakes on exactly 10\% of inputs. Then, the exact chance that is consistent on a sample is $0.9^m$, and we only need $m = 29$ to push this chance below $0.05$.

Now suppose that $|\Hy| = 2^n$. Then, $\ln |\Hy| = n \ln 2$. Thus, the sample bound grows linearly with $n$ as long as we keep $\delta$ and $\eps$ fixed.
\begin{thebibliography}{9}
\bibitem{kv}
Michael J. Kearns and Umesh V. Vazirani.
\emph{An Introduction to Computational Learning Theory}.
MIT Press, 1994. Chapter 2, ``Occam's Razor.''

\bibitem{ssbd}
Shai Shalev-Shwartz and Shai Ben-David.
\emph{Understanding Machine Learning: From Theory to Algorithms}.
Cambridge University Press, 2014. Section 2.3.1 and Corollary 2.3.
\end{thebibliography}

\end{document}
